\documentclass[10pt,a4paper]{article}

\usepackage[utf8]{inputenc}
\usepackage[T1]{fontenc}
\usepackage[spanish,es-nodecimaldot]{babel}
\usepackage[margin=1.9cm,top=2.2cm,bottom=1.9cm]{geometry}
\usepackage{booktabs}
\usepackage{array}
\usepackage{xcolor}
\usepackage{colortbl}
\usepackage{fancyhdr}
\usepackage{enumitem}
\usepackage{titlesec}
\usepackage{helvet}
\renewcommand{\familydefault}{\sfdefault}

\definecolor{critico}{RGB}{200,30,30}
\definecolor{alto}{RGB}{225,120,20}
\definecolor{medio}{RGB}{215,180,20}
\definecolor{gris}{RGB}{95,95,95}
\definecolor{grisclaro}{RGB}{240,240,240}

\titleformat{\section}{\large\bfseries\color{black}}{}{0pt}{}
\titlespacing{\section}{0pt}{11pt}{4pt}

\pagestyle{fancy}
\fancyhf{}
\renewcommand{\headrulewidth}{0.4pt}
\lhead{\small\color{gris}Reporte Ejecutivo de Gestión de Riesgos}
\rhead{\small\color{gris}Clínica --- Uso interno}
\cfoot{\small\color{gris}Página \thepage\ de 3}

\newcommand{\nivel}[2]{\textcolor{#1}{$\blacksquare$}\ #2}

\setlength{\parskip}{4pt}
\setlength{\parindent}{0pt}

\begin{document}

\begin{center}
{\LARGE\bfseries Reporte Ejecutivo de Gestión de Riesgos}\\[6pt]
{\large Registro inicial de riesgos de seguridad de la información}\\[10pt]
\end{center}

\begin{tabular}{@{}ll@{}}
\textbf{Destinatario:} & Directorio de la Clínica \\
\textbf{Emite:} & Responsable de Seguridad de la Información \\
\textbf{Alcance:} & Sistemas de información, procesos administrativos e infraestructura \\
\textbf{Elaborado con:} & SimpleRisk --- metodología Probabilidad $\times$ Impacto (escala 1--5) \\
\end{tabular}

\vspace{6pt}
\hrule
\vspace{10pt}

\section{Resumen ejecutivo}

Como resultado de la auditoría externa que identificó debilidades en la gestión
de riesgos, se conformó el primer registro formal de riesgos de seguridad de la
información de la Clínica. Se identificaron y valoraron \textbf{siete riesgos},
cada uno con su probabilidad e impacto justificados y con un plan de tratamiento
asignado.

El resultado exige atención del Directorio: \textbf{cuatro de los siete riesgos
alcanzan nivel Crítico}. Esta concentración no indica una gestión deficiente sino
la situación esperable en una organización que digitalizó su operación sin haber
constituido una función de seguridad de la información. Los riesgos existían
antes de este análisis; lo nuevo es que ahora están identificados y dimensionados.

Tres conclusiones concentran la atención del Directorio:

\begin{enumerate}[leftmargin=16pt,itemsep=3pt]
\item \textbf{La continuidad asistencial depende hoy de un único punto de falla.}
Un incidente de cifrado sobre el sistema de historias clínicas interrumpiría la
atención de 800 pacientes diarios. Las copias de seguridad existen, pero nunca se
verificó que puedan restaurarse.

\item \textbf{El riesgo más subestimado no es un ataque sino un error.} El envío
de planillas a obras sociales sin procedimiento formal alcanza el mismo nivel de
criticidad que un ataque de ransomware. La evidencia sectorial indica que la
entrega de información al destinatario equivocado es la causa más frecuente de
pérdida de datos en instituciones de salud.

\item \textbf{El plan completo requiere una inversión moderada.} Los seis planes
de acción totalizan aproximadamente \textbf{USD 32.500}, distribuidos en seis
meses. Los dos planes que atienden riesgos críticos con mayor urgencia suman
menos de USD 14.000.
\end{enumerate}

\section{Distribución de los riesgos identificados}

\begin{center}
\renewcommand{\arraystretch}{1.25}
\begin{tabular}{@{}lccl@{}}
\toprule
\textbf{Nivel} & \textbf{Cantidad} & \textbf{\% del total} & \textbf{Acción requerida} \\
\midrule
\nivel{critico}{Crítico} & 4 & 57\,\% & Acción inmediata. Escalar a Dirección. \\
\nivel{alto}{Alto}       & 2 & 29\,\% & Tratamiento prioritario a corto plazo. \\
\nivel{medio}{Medio}     & 1 & 14\,\% & Plan de acción a mediano plazo. \\
\midrule
\textbf{Total} & \textbf{7} & \textbf{100\,\%} & \\
\bottomrule
\end{tabular}
\end{center}

\newpage

\section{Top 5 riesgos por nivel}

\begin{center}
\renewcommand{\arraystretch}{1.35}
\begin{tabular}{@{}c>{\raggedright\arraybackslash}p{6.6cm}cccl@{}}
\toprule
\textbf{\#} & \textbf{Riesgo} & \textbf{Prob.} & \textbf{Imp.} & \textbf{Valor} & \textbf{Nivel} \\
\midrule
1 & Ransomware que cifra el sistema de historias clínicas & 4 & 5 & \textbf{20} & \nivel{critico}{Crítico} \\
2 & Acceso de personal a historias clínicas ajenas a su función & 4 & 4 & \textbf{16} & \nivel{critico}{Crítico} \\
3 & Compromiso de credenciales por phishing al área administrativa & 4 & 4 & \textbf{16} & \nivel{critico}{Crítico} \\
4 & Fuga de datos de obras sociales por envío inseguro de planillas & 4 & 4 & \textbf{16} & \nivel{critico}{Crítico} \\
5 & Imposibilidad de restaurar los backups ante un incidente & 3 & 5 & \textbf{15} & \nivel{alto}{Alto} \\
\bottomrule
\end{tabular}
\end{center}

\textbf{Lectura de los resultados.} Los tres primeros riesgos comprometen la
continuidad operativa o la confidencialidad de la historia clínica, que es el
activo de mayor valor y mayor sensibilidad legal de la institución. El cuarto
merece una mención específica: se trata de un error administrativo cotidiano, no
de un ataque, y alcanza el mismo nivel que los anteriores.

El quinto riesgo actúa como amplificador. Por sí solo no genera un incidente,
pero determina si el primero resulta recuperable o definitivo. Es la diferencia
entre una interrupción de horas y la pérdida irreversible de historias clínicas,
documentación cuya conservación es obligatoria.

\section{Estado de los planes de acción}

Se definieron seis planes que cubren la totalidad de los riesgos identificados.
Todos se encuentran en estado \textbf{No iniciado}, a la espera de aprobación del
Directorio.

\begin{center}
\renewcommand{\arraystretch}{1.3}
\begin{tabular}{@{}l>{\raggedright\arraybackslash}p{5.6cm}ccr@{}}
\toprule
\textbf{Plan} & \textbf{Objetivo} & \textbf{Riesgos} & \textbf{Plazo} & \textbf{Presupuesto} \\
\midrule
PA01 & Respaldo verificado y contención de propagación & R01, R05 & 90 días & USD 12.000 \\
PA02 & Trazabilidad y perfilado de accesos a la historia clínica & R02 & 60 días & USD 3.500 \\
PA03 & Segundo factor de autenticación y concientización & R03 & 120 días & USD 5.000 \\
PA04 & Procedimiento seguro de intercambio con obras sociales & R04 & 45 días & USD 1.500 \\
PA05 & Regularización del tratamiento de datos sensibles & R07 & 180 días & USD 2.500 \\
PA06 & Redundancia eléctrica y de climatización & R06 & 150 días & USD 8.000 \\
\midrule
\multicolumn{4}{r}{\textbf{Total}} & \textbf{USD 32.500} \\
\bottomrule
\end{tabular}
\end{center}

La secuencia de vencimientos no sigue el orden de severidad. PA04 vence primero
pese a no ser el riesgo de mayor valor, porque su implementación depende de
definir un procedimiento y no de adquirir tecnología. PA03 tiene el plazo más
extenso porque su componente de capacitación requiere sostenerse en el tiempo
para producir efecto verificable.

\section{Recomendaciones prioritarias}

\subsection*{1. Aprobar de inmediato los dos planes de menor costo}

PA02 y PA04 suman USD 5.000 y atienden dos riesgos de nivel Crítico. Ambos se
resuelven con controles de procedimiento antes que con inversión tecnológica, y
ofrecen la mayor reducción de riesgo por unidad de inversión de todo el plan. No
existe argumento económico para postergarlos.

\subsection*{2. Verificar la capacidad de restauración antes que cualquier otra medida}

La Clínica realiza copias de seguridad diarias, pero nunca comprobó que puedan
restaurarse. Una copia de seguridad no verificada es una suposición, no un
control. Se recomienda ejecutar una prueba de restauración completa dentro de los
próximos treinta días, con independencia del avance del resto del plan, y
documentar el tiempo de recuperación obtenido.

\subsection*{3. Designar un responsable del tratamiento de datos personales}

Los datos de salud revisten carácter de datos sensibles bajo la Ley 25.326 y
están sujetos a un régimen agravado. Actualmente la Clínica no tiene registradas
sus bases de datos, carece de política de privacidad formal y sus contratos con
proveedores no incluyen cláusulas de tratamiento de datos. La designación de un
responsable es la condición previa para regularizar esta situación.

\subsection*{4. Establecer la revisión periódica del registro de riesgos}

Este registro refleja la situación al momento de su elaboración. Se recomienda su
revisión semestral y ante todo cambio relevante de infraestructura, sistemas o
normativa. Un registro de riesgos que no se actualiza pierde validez con la misma
velocidad con que cambia la organización que describe.

\section{Consideración final}

Ninguno de los riesgos identificados es exclusivo de esta institución: responden a
patrones documentados en el sector salud a escala internacional. Lo que distingue
a una organización de otra no es la ausencia de riesgos sino la decisión de
identificarlos y tratarlos antes de que se materialicen. Este registro constituye
la respuesta formal a los hallazgos de la auditoría externa y establece la línea
base sobre la cual medir el progreso en los próximos ciclos de revisión.

\vspace{20pt}
\hrule
\vspace{8pt}
{\small\color{gris}
Elaborado con SimpleRisk. La metodología de valoración, los supuestos del
análisis y la justificación individual de cada riesgo se encuentran disponibles
en la documentación técnica que acompaña este reporte.}

\end{document}
